\chapter*{Abstract}

Modern supermarket and convenience chains in Vietnam face two tightly coupled operational bottlenecks. On the sales floor, shoppers navigating thousands of stock-keeping units (SKUs) must manually compare pack sizes, unit prices, dietary constraints, and promotions in noisy store environments while staff capacity is stretched thin. In back-office store operations, managers must make daily replenishment, markdown, supplier-return, and write-off decisions for short-shelf-life perishable goods under censored demand (where out-of-stock days hide true customer demand) and high shrinkage risk. Generic chatbots and standalone forecasting scripts fail to solve these problems reliably: unconstrained large language models (LLMs) hallucinate prices, stock availability, and unsafe health advice, incur 5--14\,s multi-call latency, and risk executing unauthorized operational actions.

To address both sides within one unified logical system, this capstone designs, builds, and evaluates a dual-agent artificial intelligence platform spanning two complementary repositories: \texttt{Retail-AI-} (a Vietnamese-first, voice-enabled sales-floor kiosk agent named \emph{Mai}, deployed live in the Singapore \texttt{sin1} serverless region) and \texttt{ops\_chest} (a conversational back-office operations agent for store managers). On the customer floor, \texttt{Retail-AI-} integrates a mobile-first Vue~3 kiosk ($\le 768$\,px frosted bottom-sheet with client-side photo/video compression $\le 640$\,px, product thumbnail cards opening a \textbf{full-screen Product Photo Lightbox} covering 99.9\% of the catalog, and real-time \textbf{Zalo Bot manager escalation}), a four-layer preloaded crossfade 2D video avatar, full-duplex three-layer anti-echo voice barge-in over local Silero VAD~v6, real-time WebSocket streaming speech-to-text (\code{/audio/transcribe/stream}) powered by Azure AI Speech SDK with Thanh Lam Mart \texttt{PhraseListGrammar} and local \texttt{DPDFNet2} ONNX neural noise suppression (zero Groq calls on live STT), progressive clause-level Azure Neural TTS with gapless audio prefetching, a \textbf{strictly on-demand mobile visual product recognition subsystem} built on a Kaggle-trained \texttt{MobileNetV3Large-224} TFLite embedder (7.04\,MB FP16 / 3.81\,MB INT8 + 1.58\,MB FP16 gallery across 3,093 SKUs, distilled via SVD from \texttt{CLIP-ViT-B/32} and \texttt{DINOv2-Small}) paired with \code{libjpeg} DCT preprocessing and a sibling-pack-size-guarded Hybrid VLM route, deterministic sub-3\,ms pre-LLM safety gates, a \textbf{multi-source three-tier product knowledge store cleaned through an 8-stage zero-defect pipeline}, and a \textbf{tiered hybrid streaming architecture} uniting a zero-planner \textbf{free Groq LPU fast-lane} (\code{qwen/qwen3.8-27b} rotating 4 keys), a zero-token Turn-1 semantic cache (\code{consultation_cache.py}), and structured JSON reasoning on DeepSeek (\code{deepseek-v4.1-flash:free} $\to$ \code{deepseek-v4.1-flash-1017}). On the operations side, \texttt{ops\_chest} combines hourly censored-demand recovery, 29-feature LightGBM quantile demand forecasting (P10/P50/P90), a profit-gated ML-versus-EWMA forecast router, perishable Newsvendor and First-Expired-First-Out (FEFO) lot-level replenishment simulation, 31 read/proposal agent tools with zero direct-execution rights, a runtime numeric grounding and false-execution guard, and strict Human-in-the-Loop (HITL) manager approval.

Development and evaluation drew on a 3,099-SKU Thanh Lam Mart catalog snapshot (3,081 validated after schema quarantine of 18 price anomalies across 699 brands; \textbf{2,137/3,099 SKUs (69.0\%)} enriched with clean product knowledge including \textbf{1,214 + 18} verified ingredient records, \textbf{1,064+} quantitative/INCI back-of-pack label lists, and \textbf{0 residual defects} after eliminating 4,606 raw defects across 11 QA dimensions), a 100-question SOTA~V2 retail/adversarial benchmark (113 turns), an independent 100-question 5-framework expert catalog audit, a 5-engine Vietnamese speech-in-noise benchmark (clean, 12\,dB, and 5\,dB supermarket noise), live HTTPS latency measurements on Vercel Singapore (\texttt{sin1}), and FreshRetailNet-50K real perishable demand across 50 held-out stores (2,675 SKU-store time series; 164 fresh SKUs over 90 days in Demo Store~18). Quantitative results show that on \texttt{Retail-AI-}, systemic engineering raised strict SOTA~V2 accuracy from 59/100 to 94/100 (96/100 functional) with 0\% XSTest over-refusal and 10/10 red-team defense; tiered \textbf{Groq LPU streaming + Turn-1 semantic caching + Singapore \texttt{sin1} edge co-location} slashed live production HTTPS latency from 2.4--5.5\,s down to \textbf{508--870\,ms} cold (TTFT \textbf{319--572\,ms} at \$0; dinner basket $<150\text{k}$ dropped from 5.50\,s to \textbf{0.87\,s}), \textbf{125--140\,ms} HTTPS on semantic paraphrase cache hits (11--36\,ms local, 0 tokens), and \textbf{90--104\,ms} session creation; migrating from post-utterance WAV upload to live WebSocket streaming STT cut post-speech wait from 1,937--3,042\,ms to 257--558\,ms (up to $9.7\times$ faster) with 0.0\% brand WER under 12\,dB supermarket noise; meanwhile, the strictly on-demand \texttt{MobileNetV3Large-224} TFLite mobile recognizer achieved \textbf{92.67\% Top-1 / 99.33\% Top-5} on a 150-SKU 3D multi-angle stress test (5.49\,ms FP16 core + 0.07\,ms 3,093-SKU search), \textbf{23.4--27.4\,ms} live HTTP warm Fast-Path latency ($\approx 150\times$ faster than cloud VLM), and \textbf{93.8\% Brand@1} with \textbf{100\%} out-of-distribution rejection in Hybrid mode on a zero-leakage smartphone evaluation set (16 unseen OpenFoodFacts photos + 4 real OOD negatives). On \texttt{ops\_chest}, hourly stockout recovery lowered demand WAPE from 32.3\% to 20.5\% (bias $-32.3\% \to -0.1\%$ over $n=200{,}000$ days), LightGBM~v3 reduced evaluation WAPE from 35.1\% to 30.2\% with 81\% interval coverage, and the out-of-sample profit-gated router (\texttt{B-Gate}) raised median store gross profit by $+4.9\%$ over the heuristic baseline ($+0.8\%$ over EWMA, winning on 76\% of held-out stores) while reducing waste to 4.8\% and cutting traffic-driver stockouts from 8.7\% to 3.4\%. Automated verification confirms 2,900+ backend tests (562 latency/cache/streaming tests, 146 knowledge/support/demo-sales tests, and 0 Ruff errors), 150+ frontend Vitest tests (25/25 in \code{App.test.ts} including the product photo modal), 20 \texttt{ops\_chest} test suites, and 20/20 runtime numeric grounding. Key remaining limitations---unbridged real-time POS synchronization between repositories, remaining unlabelled nutrition/calorie SKUs, and absence of multi-week physical supermarket floor trials---are explicitly scoped as prioritized future work.

\vspace{0.6em}
\noindent\textbf{Keywords:} Retail AI agents; agentic commerce; real-time streaming STT; neural speech enhancement; Groq LPU SSE streaming; semantic caching; multi-source data cleaning; perishable inventory optimization; LightGBM quantile forecasting; human-in-the-loop governance; Vietnam retail.
